\clearpage
\begin{center}
{\Large\bfseries EXPOSÉ XXII}\par
\vspace{1.5cm}
{\Large\bfseries A CONGRUENCE FORMULA FOR THE ZETA FUNCTION}\par
\vspace{0.8cm}
{\large by N. Katz}\par
\end{center}

\vspace{1.5cm}
\section*{0. Introduction}

Let $X$ be a proper scheme over $\mathbb F_q$. Its zeta function, regarded as
an element of $\mathbb Z[[T]]$, may be reduced modulo $p$; we wish to compute
it modulo $p$. We obtain an expression of the kind one might have expected
from fixed-point formulas of ``Woods Hole'' type: an alternating product of
characteristic polynomials of a suitable ``Frobenius'' acting on any one of
the following three cohomology spaces:

\begin{enumerate}[label=\arabic*)]
\item coherent cohomology $H^i(X,\mathcal O_X)$ (3.1.1);
\item de Rham cohomology
  $H^i_{\mathrm{DR}}(X)=\mathbb H^i(X,\Omega_X^\bullet)$ (3.1.1);
\item étale cohomology $H^i_{\mathrm{et}}(X,\mathbb Z/p\mathbb Z)$ (3.1.2).
\end{enumerate}

This expression should be viewed as an improvement of a special case of the
Woods Hole trace formula, namely that of the structure sheaf and the
Frobenius endomorphism. That case is indeed a corollary of the congruence
formula (cf. 3.2.1).

The method of this exposé is completely naïve. It consists in reducing to the
case of a hypersurface and then ``reducing modulo $p$'' Dwork's computation
of the actual zeta function of a hypersurface. By ``reducing'' carefully, one
obtains the improvement 5.0.6, due to Ax~[1] (and also to Dwork in the smooth
case [4, \S 7]), of the well-known Warning--Chevalley theorem~[8].

Another method, entirely cohomological and in the spirit of SGA~5, which also
applies to ``twisted coefficients,'' has just been found by Deligne and ought
to appear in the new edition of SGA~5. It rests on the trace formula of Woods
Hole type and on Deligne's ``symmetric Künneth formula''
(SGA~4, XVII~5.5.21).

\section*{1. A reminder on ``$p$-linear'' operations}

\subsection*{1.0.}
Let $V$ be a finite-dimensional vector space over a field $k$ of
characteristic $p>0$, and let $q$ be a power of $p$. A map

\begin{equation}\tag{1.0.1}
\varphi:V\longrightarrow V
\end{equation}

is called $q$-\emph{linear} if it is additive and satisfies

\begin{equation}\tag{1.0.2}
\varphi(\lambda v)=\lambda^q\varphi(v)
\qquad\text{for every }\lambda\in k,\ v\in V.
\end{equation}

\subsection*{1.0.3.}
Notice that the $n$-th iterate $\varphi^n$ of a $q$-linear endomorphism is
$q^n$-linear and that, if $k$ is not perfect, the image $\varphi(V)$ of a
$q$-linear endomorphism need not be a $k$-vector subspace.

\subsection*{1.0.4.}
For every extension $K/k$, a $q$-linear endomorphism $\varphi$ of $V$ gives
rise to a unique $q$-linear endomorphism $\varphi_K$ of
$V_K=V\otimes_kK$ satisfying

\begin{equation}\tag{1.0.5}
\varphi_K(\mu\otimes v)=\mu^q\varphi(v)
\qquad\text{for }\mu\in K,\ v\in V.
\end{equation}

\subsection*{1.0.6.}
A $q$-linear endomorphism $\varphi$ of $V$ is called \emph{semisimple} if
its image spans $V$ over $k$, and \emph{nilpotent} if some iterate is zero.
Define the following $\varphi$-stable $k$-subspaces:

\begin{equation}\tag{1.0.7}
V_{\mathrm{ss}}=
\bigcap_{n\geq1}\langle\varphi^n(V)\rangle_k,
\end{equation}

\begin{equation}\tag{1.0.8}
V_{\mathrm{nilp}}=
\bigcup_{n\geq1}\Ker(\varphi^n:V\longrightarrow V).
\end{equation}

Clearly, the restriction of $\varphi$ to $V_{\mathrm{ss}}$ is semisimple,
while its restriction to $V_{\mathrm{nilp}}$ is nilpotent. If $k$ is
perfect, then

\begin{equation}\tag{1.0.9}
V=V_{\mathrm{ss}}\oplus V_{\mathrm{nilp}}.
\end{equation}

If $k$ is not perfect, however, (1.0.9) is false, since there exist injective
$q$-linear endomorphisms which are not semisimple. For an arbitrary extension
$K/k$, one has

\begin{equation}\tag{1.0.10}
(V_K)_{\mathrm{ss}}=V_{\mathrm{ss}}\otimes_kK.
\end{equation}

